\author{Schmid Lothar}
\documentclass[11pt]{mybib}

\setheaderlogo{mnr-green}


% Absatzformatierung
% \setlength{\parindent}{0pt}
% \setlength{\parskip}{0.6em}
% \renewcommand{\arraystretch}{1.8}
\begin{document}
    \section{2. Petrus 3,8-9}
    \subsection{Bibeltext}
    \begin{auslegung}{2. Petrus 3.8-9}    
        \konj[b, p]{Dieses} eine aber sollt \person[b, p]{ihr} nicht \verbN[b, p]{übersehen}, \person[b, p]{Geliebte},
        \konj[b, p]{dass} ein Tag bei dem \person[b, p]{Herrn} \verbN[b, p]{ist} wie tausend Jahre,
        \konj[b, p]{und} tausend Jahre wie ein Tag!   
        \switchcolumn
        Mir \betonung[b, p]{Dieses} will Petrus die nachfolgende Aussage bekräftigen. \betonung[b, p]{Dieses} genau das, was er jetzt sagen wird, müssen sie beachten.
        \switchcolumn*
        Der Herr zögert nicht die Verheissung hinaus, wie etliche es für ein Hinauszögern halten, sondern er ist langmütig gegen uns, weil er nicht will, dass jemand verloren gehe, sondern dass jedermann Raum zur Busse habe.
        \switchcolumn
        Hier die Erklärung
    \end{auslegung}
    
\end{document}